The original study from \cite{ORIGINAL_STUDY} addresses the problem of ramp
metering for a bottleneck located far downstream of the metered ramp. Conventional
linear feedback-type controllers, such as ALINEA \citep{ALINEA} or PI-ALINEA \citep{PI-ALINEA},
suffer from stability issues in this setting due to the significant time-delay between
the ramp and the distant bottleneck: by the time metered vehicles reach the bottleneck,
its traffic state may have changed substantially from when it was sampled to compute
the metering rate. Rather than compensating for this delay through predictors of traffic
flow evolution, the authors propose formulating the problem as a Q-learning problem in
which an agent learns a non-linear feedback policy directly from interaction with a
traffic simulation environment.

The state vector consists of four continuous variables: traffic densities measured at
three representative locations along the stretch between the ramp and the downstream
bottleneck ($\rho_1$, $\rho_2$, $\rho_3$), and an estimated ramp demand $D_\text{ramp}$.
The agent selects a discrete metering rate at each time step, with the action space
bounded above by the estimated ramp demand to avoid unrealistic control actions. The
reward function penalises deviations of the downstream bottleneck density $\rho_3$ from
a desired critical density value.

To handle the computational burden of the four-dimensional continuous state space, the
Q-value function is approximated by an artificial neural network (ANN) rather than a
lookup table, reducing the number of parameters to learn from an estimated 25.6 million
Q-values to approximately 67,000. State variables are first encoded via a tile coding
scheme into a 140-dimensional binary feature vector before being passed to the ANN.
Parameters are updated online at each time step using an incremental back-propagation
algorithm \citep{Backpropagation}.

The approach is validated on a synthetic freeway section implemented in the Cell
Transmission Model (CTM), consisting of a three-lane mainline with a lane-drop to two
lanes located 3500\,m downstream of the metered ramp.
Learning is reported to have converged after
approximately 0.7 million episodic iterations. The proposed method is compared against
an uncontrolled baseline and a conventional PI controller (PI-ALINEA \citep{PI-ALINEA}). Results show that the learned
policy successfully maintains the bottleneck density close to the desired value with
negligible oscillations, whereas the PI controller exhibits some instability at this
distance. A robustness analysis under additive white noise on traffic demands shows
satisfactory performance up to a standard deviation of 200\,veh/h.
%\clearpage


